\clearpage
\section{Conclusions}
\label{sec:conclusions-final}
\begingroup
\fontsize{10.2}{12.8}\selectfont
\setlength{\parskip}{3pt plus .5pt}

The geometric thesis is established through reconstruction and transport:
a positive readout of the discrete structure of the continuum preserves,
with its markings, the record from which it arises and the readers that
record determines. Total charge and ratios of minors recover each
coordinate; the two nonadic memories provide another operator-theoretic
reconstruction. Value and form are thus articulated by explicit inverses,
after APP--TRIT--TPK generation and its compatible prolongations.
The double projection keeps the arithmetic realization and
exceptional incidence together within that genealogy.

Refinement overcomes the limitation of a matrix with thirteen columns.
For any finite rank and external dimension, a sufficiently high nonadic
level produces the positive minors and the corresponding amplituhedral
image. Inherited nodes preserve their flags; cell
averages preserve the measure and record the orthogonal
detail. In rank one, simplicial coverage and residue recursion
extend to every finite dimension. In higher ranks,
the result determines compatible positive collections with proved rank and
limit; their dimensional scope is distinguished from the
global classification of all amplituhedral cells.

Conservation admits three complementary expressions. The canonical
form adds under subdivision of a fixed support; effective
energy is preserved by Schur reduction; and the chronological
measure determines moments whose increments decompose into
recoverable details. These moments construct Jacobi
operators whose finite compressions have simple spectrum and positive residues. This yields
an effective passage from arithmetic separations to a
spectral representation, with a continued fraction determined by the
same measure. Mutations changing only the chart
preserve the readers of the marked point.

The physical realization gains precision through its operators. The
Hamiltonians of the three tritic regimes have explicit
symplectic forms and Legendre transforms. The variable
elliptic chart incorporates its connection term; parabolic
memory preserves its invariant; hyperbolic opening has
its own evolution. The readers of action and constitutive
response are transported through the Grassmannian inverse. With
the declared regional outputs and integer markings,
the labelled constitutive response in turn recovers the record.
The equality-of-traces condition governs the gluing of forms
valued in the action line.

The exceptional deformation preserves volume and ternary
symmetry, and converts parameter inversion into lattice
duality. Its theta reciprocity and split-signature realization
link geometric variation with lattice duality.
In Yang--Mills, the areal--volumetric scale, common measure,
uniform coercivity and persistent witness place the gap
alongside the Hamiltonian realizing it. The development preserves
the physical space and the intertwiners required to transport
that assertion through refinement.

The joint result is an architecture of forms with memory:
dimension can increase, representation can change and
the relevant information remains recoverable through
specified maps. This is the article's distinctive content.
Positive geometry, lattice duality and the spectral
realization are linked by constructions specifying
what they generate, what they preserve and how they are reconstructed
from the same discrete origin.
The incidence representation also preserves the uniform component
of the record: its inverse reconstructs all twelve coordinates and allows
the geometric direction to be recovered after memory transport.
In the sectorial realization, total area and weighted occupation
jointly determine the hexad--octad distribution. The specified
purification converts that distribution into an entanglement structure,
with a finite identity relating entropy, sector areas and relative entropy.
Information conservation is thus specified for each application:
the full record, direction, area and composition have explicitly
distinct inverses or fibres.
The area scale is determined by a length generated through the linear
composition of recorded return and the local Hadamard norm.
Radial reciprocity determines the gravitational coefficient in the
declared reduced realization and subsequently identifies that length as
the Planck length. Mutations of the same marked point and
inherited refinements preserve the area operator, its
multiplicities and the modular identity. Canonical forms
valued in that operator extend conservation to residues,
while each horizon radius retains its own geometric
factor. This composition makes the dependence between
length, incidence and area explicit without altering the sectorial states.
\par\endgroup
\clearpage
